\section{Prompting Without a Structural Approach}

Development in the first part had no predefined process for how prompts should be written or sequenced. The agent --- used via Claude Code's chat sidebar in Visual Studio Code (Section 3.2) --- was simply given a feature to build and prompted feature by feature, with the structure of each prompt worked out ad hoc rather than following a fixed template. The agent operated throughout under a persistent project file, CLAUDE.md, which supplied backend and frontend conventions and a set of architectural rules and is covered separately in Section 4.3. What this file did not provide was a per-feature planning or review step: no design was drafted or approved for a given feature before implementation began, and the file itself was updated reactively over the course of the project, as problems surfaced, rather than being fixed in advance and left unchanged. \Cref{lst:login-prompt} shows a representative example of this ad hoc structure, taken from the prompt that initiated the login component.

\begin{listing}[H]
\begin{center}
\begin{promptbox}
I created a react app in the frontend folder, i want to create the first login components, i want you to develop this components for me. I already created a Mockup for the component and i want you to make it look exactly the same as the mockup, so if you want to ask me something about the mockup details it's okay. But don't change in the code until you ask me.
@frontend/api-craft-app/src/components/login.js
\end{promptbox}
\end{center}
\caption{Example prompt used to initiate development of the login component.}
\label{lst:login-prompt}
\end{listing}

A further small set of rules was established early and carried through the rest of Part 1, layered on top of the project file and applied at the level of individual prompts: the agent was told to ask clarifying questions before generating code where the specification was unclear, and it was not permitted to modify any existing file without asking for permission first, so that every change could be reviewed before being applied. Later prompts added a further rule requiring the agent to briefly explain what it proposed to change and why before applying it. These rules shaped how the agent behaved within a session, but they did not amount to a planning step performed before a feature was started --- there was no point at which a design was reviewed and approved ahead of writing code.

The typical pattern for a new frontend feature began with a mockup, sometimes accompanied by a short verbal description of details such as colors, fonts, and layout, followed by a single prompt asking the agent to reproduce it. The result was then compared against the mockup and refined through further prompts, usually several rounds of them, before the layout was considered close enough to the intended design. This was true even when the mockup was provided upfront: matching it exactly in a single pass did not happen in practice, and the gap was closed through iterative correction afterward rather than through a more complete initial specification.

Backend prompts followed a different pattern, since there was no visual reference to match against. Instead of iterating toward a mockup, a single prompt typically specified the data model or endpoint involved, the validation rules to apply, and the expected behavior together in one request, often alongside explicit constraints on scope and performance. This combined pattern --- data model or endpoint, validation, and constraints specified together in a single prompt --- recurred across backend features throughout Part 1, in contrast to the frontend's iterative, comparison-driven refinement against a visual target. \Cref{lst:auth-prompt} shows a representative example of this pattern, taken from the prompt that initiated the authentication feature.

\begin{listing}[H]
\begin{center}
\begin{promptbox}
I want you to implement a basic user authentication flow including both backend and minimal frontend integration. First, create a User model in the backend with the fields id (primary key), username (unique and required), and password (required), and store it in a database table. Then, implement a signup endpoint (POST /api/auth/signup) that accepts a username and password, validates that both are provided, checks whether the username already exists, and if not, creates and saves a new user in the database and returns a success response. Next, implement a login endpoint (POST /api/auth/login) that accepts username and password, validates them against the stored data, and returns a success response if correct or an error otherwise. On the frontend side, after a successful signup, the user should automatically be redirected to the login page, and after a successful login, the user should be redirected to a simple page displaying "Login successful." Keep the implementation simple without adding extra features such as JWT, sessions, roles, or OAuth, and do not introduce unnecessary complexity. Only implement the required logic and endpoints, and do not modify any existing files without asking for my permission.
\end{promptbox}
\end{center}
\caption{Example prompt used to initiate development of the authentication feature.}
\label{lst:auth-prompt}
\end{listing}

Over the course of Part 1, individual task prompts grew progressively more detailed, but this growth happened reactively, in response to problems encountered in earlier features, rather than as part of a planned prompting strategy. Constraints that appear in later prompts --- for instance, explicit instructions to avoid unnecessary re-renders, avoid repeated API calls, or not place API calls inside render logic --- do not appear in the earliest prompts of the project. The same pattern holds for scope constraints such as "do not modify unrelated files" or "do not add features beyond what is requested," which became a fixed part of later prompts after earlier, looser prompts had produced unwanted side effects. The prompting approach in Part 1 can therefore be described as incrementally, rather than deliberately, structured at the task level: the rules that accumulated in individual prompts came out of correcting specific problems as they appeared, not from a plan set before development began --- and the same was true one level up, where the persistent project file supplying broader conventions started out with only a partial set of them and grew the rest the same reactive way.

